\documentclass[11pt,a4paper]{article}
\usepackage[utf8]{inputenc}
\usepackage[T1]{fontenc}
\usepackage[catalan]{babel}
\usepackage{lmodern}
\usepackage[sfdefault,lining]{FiraSans}
\usepackage[a4paper,top=2.3cm,bottom=2.2cm,left=1.8cm,right=1.8cm,headheight=1.3cm,headsep=0.4cm,footskip=1cm]{geometry}
\usepackage{amsmath,amssymb}
\usepackage{icomma}
\usepackage{graphicx}
\usepackage[table]{xcolor}
\usepackage{tikz}
\usepackage{fontawesome5}
\usepackage{ccicons}
\usepackage{fancyhdr}
\usepackage{lastpage}
\usepackage{tasks}
\usepackage{needspace}
\usepackage[most]{tcolorbox}
\usepackage{array}
\usepackage[hidelinks,pdftitle={Solucions de la fitxa de la unitat 6: Vectors i rectes. 4t d'ESO Matemàtiques B},pdfsubject={Matemàtiques B 4t ESO},pdfkeywords={vectors, components, mòdul, producte escalar, punt mig, distància, punts alineats, equacions de la recta, pendent, posició relativa, rectes perpendiculars, mediatriu}]{hyperref}

\definecolor{navy}{HTML}{1B2748}
\definecolor{accent}{HTML}{E8603A}
\definecolor{teal}{HTML}{1F9E90}
\definecolor{soft}{HTML}{EEF1F8}
\definecolor{muted}{HTML}{5D6685}

\setlength{\parindent}{0pt}
\setlength{\parskip}{3pt}
\settasks{label=\alph*), label-format=\bfseries\color{navy}, label-width=1.3em, item-indent=1.8em, column-sep=0.6em, after-item-skip=5pt, before-skip=3pt, after-skip=4pt}

% ---------- capçalera i peu ----------
\newcommand{\llicencia}{\href{https://creativecommons.org/licenses/by-nc-sa/4.0/deed.ca}{\ccbyncsa}\ \ Llicència Creative Commons Reconeixement-NoComercial-CompartirIgual 4.0 Internacional (CC BY-NC-SA 4.0)}
\pagestyle{fancy}
\fancyhf{}
\fancyhead[L]{\raisebox{-0.35\height}{\includegraphics[height=0.95cm]{logo-gva.png}}}
\fancyhead[C]{\small\color{muted}Solucions de la fitxa · Unitat 6: Vectors i rectes · 4t ESO B}
\fancyhead[R]{\raisebox{-0.35\height}{\includegraphics[height=0.95cm]{logo-vermellar.png}}}
\fancyfoot[L]{\scriptsize\color{muted}\llicencia}
\fancyfoot[R]{\small\color{muted}\thepage}
\renewcommand{\headrulewidth}{0.4pt}
\renewcommand{\headrule}{\color{muted!40}\hrule height \headrulewidth}
\fancypagestyle{primera}{\fancyhf{}\renewcommand{\headrulewidth}{0pt}
  \fancyfoot[L]{\scriptsize\color{muted}\llicencia}
  \fancyfoot[R]{\small\color{muted}\thepage}}

% ---------- nivells i exercicis ----------
\newcommand{\estrelles}[1]{{\color{accent}\ifcase#1\or\faStar\faStar[regular]\faStar[regular]\or\faStar\faStar\faStar[regular]\or\faStar\faStar\faStar\fi}}
\newcounter{exer}
\newcommand{\ex}[1]{\par\needspace{4\baselineskip}\medskip\refstepcounter{exer}\noindent{\color{navy}\bfseries\theexer.}\ \ {\small\estrelles{#1}}\ \ }

\newcommand{\seccio}[2]{\par\needspace{16\baselineskip}\bigskip
  \begin{tcolorbox}[colback=navy,colframe=navy,coltext=white,arc=2mm,boxrule=0pt,left=3mm,right=3mm,top=1.2mm,bottom=1.2mm]
  {\large\bfseries #1\quad #2}
  \end{tcolorbox}\vspace{-2pt}}

\newcommand{\exf}[1]{\par\medskip\refstepcounter{exer}\noindent{\color{navy}\bfseries\theexer.}\ \ {\small\estrelles{#1}}\ \ }
\newcolumntype{C}[1]{>{\centering\arraybackslash}p{#1}}
% quadre de recordatori (només a la fitxa d'exercicis)
\newcommand{\recorda}[1]{\par\smallskip\begin{tcolorbox}[colback=accent!7,colframe=accent!55,arc=2mm,boxrule=0.5pt,left=2.5mm,right=2.5mm,top=1mm,bottom=1mm]\small{\color{accent}\faLightbulb[regular]}\ \textbf{Recorda:}\ #1\end{tcolorbox}}
\usetikzlibrary{arrows.meta}
\usepackage{multicol}
\usepackage{enumitem}
\definecolor{violet}{HTML}{7C3AED}
\definecolor{orange}{HTML}{EA8A1F}
\newcommand{\res}[1]{\boxed{#1}}\newcommand{\rest}[1]{\fbox{#1}}

\def\fmneed{11}% que el títol d'una secció no quede sol al final del full
% =====================================================================
% Capçalera FentMates de les fitxes: un color per a cada curs.
% Els colors són els de la pàgina inicial (LEVELS a assets/js/materials.js);
% el logo ve de _eines/logo/pdf/ (python3 _eines/logo/genera.py).
%
% Ús, al final del preàmbul de la unitat:
%   % =====================================================================
% Capçalera FentMates de les fitxes: un color per a cada curs.
% Els colors són els de la pàgina inicial (LEVELS a assets/js/materials.js);
% el logo ve de _eines/logo/pdf/ (python3 _eines/logo/genera.py).
%
% Ús, al final del preàmbul de la unitat:
%   % =====================================================================
% Capçalera FentMates de les fitxes: un color per a cada curs.
% Els colors són els de la pàgina inicial (LEVELS a assets/js/materials.js);
% el logo ve de _eines/logo/pdf/ (python3 _eines/logo/genera.py).
%
% Ús, al final del preàmbul de la unitat:
%   \input{../comu/fentmates.tex}
%   \fmcurs{1eso}
% i, al començament del document:
%   \capcalera{Matemàtiques}{1r ESO}{Fitxa d'exercicis · Unitat 2}{Divisibilitat}
% Opcional abans de \input: \def\fmneed{14} (espai mínim abans de cada secció).
% =====================================================================
\makeatletter
\def\fm@logodir{../../logo/pdf}
\providecommand{\fmneed}{10}
\def\fm@cursnom{}

% clau del curs, color i color del text damunt del color
\newcommand{\fm@curs}[3]{\@namedef{fm@color@#1}{#2}\@namedef{fm@text@#1}{#3}}
\fm@curs{1eso}{E6B84A}{navy}
\fm@curs{2eso}{3B82F6}{white}
\fm@curs{3eso}{FF7A50}{white}
\fm@curs{4eso-a}{A78BFA}{white}
\fm@curs{4eso-b}{6D28D9}{white}
\fm@curs{1batx-mi}{9B1C1C}{white}
\fm@curs{1batx-ccss}{EF6B6B}{white}
\fm@curs{2batx-mii}{1E6B47}{white}
\fm@curs{2batx-ccss}{2BB3A3}{white}

\newcommand{\fmcurs}[1]{%
  \@ifundefined{fm@color@#1}{\PackageError{fentmates}{Curs desconegut: #1}{}}{}%
  \edef\fm@tmp{\noexpand\definecolor{curs}{HTML}{\@nameuse{fm@color@#1}}%
               \noexpand\colorlet{curstext}{\@nameuse{fm@text@#1}}}\fm@tmp
  \colorlet{cursfosc}{curs!72!black}%
  \def\fm@logo{\fm@logodir/vertical-#1.pdf}%
  \def\fm@marca{\fm@logodir/marca-#1.pdf}}

% ---------- capçalera de la primera pàgina ----------
% #1 matèria (petit), #2 curs, #3 tipus de document i unitat, #4 títol
\newcommand{\capcalera}[4]{%
  \gdef\fm@cursnom{#2}%
  \par\noindent
  \begin{tcolorbox}[colback=curs,colframe=curs,boxrule=0pt,arc=2.2mm,width=3.5cm,height=2.15cm,
      valign=center,left=3mm,right=2mm,top=1mm,bottom=1mm,nobeforeafter,box align=center]
    \color{curstext}{\footnotesize #1}\par\vspace{1pt}{\fontsize{19}{21}\selectfont\bfseries #2\par}%
  \end{tcolorbox}\hspace{5mm}%
  \begin{minipage}[c]{\dimexpr\linewidth-3.5cm-5mm-2.9cm\relax}
    {\color{cursfosc}\small\bfseries\MakeUppercase{#3}}\par\vspace{2pt}
    {\color{navy}\fontsize{20}{23}\selectfont\bfseries #4\par}
  \end{minipage}\hfill
  \begin{minipage}[c]{2.6cm}\raggedleft\includegraphics[width=2.6cm]{\fm@logo}\end{minipage}\par
  \vspace{2pt}}

% ---------- barres de secció del color del curs ----------
\renewcommand{\seccio}[2]{\par\needspace{\fmneed\baselineskip}\bigskip
  \begin{tcolorbox}[colback=curs,colframe=curs,coltext=curstext,arc=2mm,boxrule=0pt,left=3mm,right=3mm,top=1.2mm,bottom=1.2mm]
  {\large\bfseries #1\quad #2}
  \end{tcolorbox}\vspace{-2pt}}

% ---------- estrelles de dificultat del color del curs (a Batxillerat no se'n posen) ----------
\@ifundefined{estrelles}{}{%
  \renewcommand{\estrelles}[1]{{\color{curs}\ifcase#1\or\faStar\faStar[regular]\faStar[regular]\or\faStar\faStar\faStar[regular]\or\faStar\faStar\faStar\fi}}}

% ---------- capçalera i peu de les pàgines ----------
% primera pàgina: logos institucionals al peu
\fancypagestyle{primera}{\fancyhf{}\renewcommand{\headrulewidth}{0pt}
  \fancyfoot[L]{\raisebox{-0.32\height}{\includegraphics[height=0.78cm]{logo-gva.png}}\hspace{5mm}%
    \raisebox{-0.32\height}{\includegraphics[height=0.78cm]{logo-vermellar.png}}}
  \fancyfoot[R]{\scriptsize\color{muted}\href{https://creativecommons.org/licenses/by-nc-sa/4.0/deed.ca}{\ccbyncsa}\ \ CC BY-NC-SA 4.0 · fentmates.org\qquad\small\thepage}}
% la resta: el símbol a l'esquerra i el curs a la dreta, amb el títol al centre com abans
\fancyhead[L]{\raisebox{-0.3\height}{\includegraphics[height=0.62cm]{\fm@marca}}}
\fancyhead[R]{\small\bfseries\color{cursfosc}\fm@cursnom}
\renewcommand{\headrule}{\color{curs!45}\hrule height \headrulewidth}
\makeatother

%   \fmcurs{1eso}
% i, al començament del document:
%   \capcalera{Matemàtiques}{1r ESO}{Fitxa d'exercicis · Unitat 2}{Divisibilitat}
% Opcional abans de \input: \def\fmneed{14} (espai mínim abans de cada secció).
% =====================================================================
\makeatletter
\def\fm@logodir{../../logo/pdf}
\providecommand{\fmneed}{10}
\def\fm@cursnom{}

% clau del curs, color i color del text damunt del color
\newcommand{\fm@curs}[3]{\@namedef{fm@color@#1}{#2}\@namedef{fm@text@#1}{#3}}
\fm@curs{1eso}{E6B84A}{navy}
\fm@curs{2eso}{3B82F6}{white}
\fm@curs{3eso}{FF7A50}{white}
\fm@curs{4eso-a}{A78BFA}{white}
\fm@curs{4eso-b}{6D28D9}{white}
\fm@curs{1batx-mi}{9B1C1C}{white}
\fm@curs{1batx-ccss}{EF6B6B}{white}
\fm@curs{2batx-mii}{1E6B47}{white}
\fm@curs{2batx-ccss}{2BB3A3}{white}

\newcommand{\fmcurs}[1]{%
  \@ifundefined{fm@color@#1}{\PackageError{fentmates}{Curs desconegut: #1}{}}{}%
  \edef\fm@tmp{\noexpand\definecolor{curs}{HTML}{\@nameuse{fm@color@#1}}%
               \noexpand\colorlet{curstext}{\@nameuse{fm@text@#1}}}\fm@tmp
  \colorlet{cursfosc}{curs!72!black}%
  \def\fm@logo{\fm@logodir/vertical-#1.pdf}%
  \def\fm@marca{\fm@logodir/marca-#1.pdf}}

% ---------- capçalera de la primera pàgina ----------
% #1 matèria (petit), #2 curs, #3 tipus de document i unitat, #4 títol
\newcommand{\capcalera}[4]{%
  \gdef\fm@cursnom{#2}%
  \par\noindent
  \begin{tcolorbox}[colback=curs,colframe=curs,boxrule=0pt,arc=2.2mm,width=3.5cm,height=2.15cm,
      valign=center,left=3mm,right=2mm,top=1mm,bottom=1mm,nobeforeafter,box align=center]
    \color{curstext}{\footnotesize #1}\par\vspace{1pt}{\fontsize{19}{21}\selectfont\bfseries #2\par}%
  \end{tcolorbox}\hspace{5mm}%
  \begin{minipage}[c]{\dimexpr\linewidth-3.5cm-5mm-2.9cm\relax}
    {\color{cursfosc}\small\bfseries\MakeUppercase{#3}}\par\vspace{2pt}
    {\color{navy}\fontsize{20}{23}\selectfont\bfseries #4\par}
  \end{minipage}\hfill
  \begin{minipage}[c]{2.6cm}\raggedleft\includegraphics[width=2.6cm]{\fm@logo}\end{minipage}\par
  \vspace{2pt}}

% ---------- barres de secció del color del curs ----------
\renewcommand{\seccio}[2]{\par\needspace{\fmneed\baselineskip}\bigskip
  \begin{tcolorbox}[colback=curs,colframe=curs,coltext=curstext,arc=2mm,boxrule=0pt,left=3mm,right=3mm,top=1.2mm,bottom=1.2mm]
  {\large\bfseries #1\quad #2}
  \end{tcolorbox}\vspace{-2pt}}

% ---------- estrelles de dificultat del color del curs (a Batxillerat no se'n posen) ----------
\@ifundefined{estrelles}{}{%
  \renewcommand{\estrelles}[1]{{\color{curs}\ifcase#1\or\faStar\faStar[regular]\faStar[regular]\or\faStar\faStar\faStar[regular]\or\faStar\faStar\faStar\fi}}}

% ---------- capçalera i peu de les pàgines ----------
% primera pàgina: logos institucionals al peu
\fancypagestyle{primera}{\fancyhf{}\renewcommand{\headrulewidth}{0pt}
  \fancyfoot[L]{\raisebox{-0.32\height}{\includegraphics[height=0.78cm]{logo-gva.png}}\hspace{5mm}%
    \raisebox{-0.32\height}{\includegraphics[height=0.78cm]{logo-vermellar.png}}}
  \fancyfoot[R]{\scriptsize\color{muted}\href{https://creativecommons.org/licenses/by-nc-sa/4.0/deed.ca}{\ccbyncsa}\ \ CC BY-NC-SA 4.0 · fentmates.org\qquad\small\thepage}}
% la resta: el símbol a l'esquerra i el curs a la dreta, amb el títol al centre com abans
\fancyhead[L]{\raisebox{-0.3\height}{\includegraphics[height=0.62cm]{\fm@marca}}}
\fancyhead[R]{\small\bfseries\color{cursfosc}\fm@cursnom}
\renewcommand{\headrule}{\color{curs!45}\hrule height \headrulewidth}
\makeatother

%   \fmcurs{1eso}
% i, al començament del document:
%   \capcalera{Matemàtiques}{1r ESO}{Fitxa d'exercicis · Unitat 2}{Divisibilitat}
% Opcional abans de \input: \def\fmneed{14} (espai mínim abans de cada secció).
% =====================================================================
\makeatletter
\def\fm@logodir{../../logo/pdf}
\providecommand{\fmneed}{10}
\def\fm@cursnom{}

% clau del curs, color i color del text damunt del color
\newcommand{\fm@curs}[3]{\@namedef{fm@color@#1}{#2}\@namedef{fm@text@#1}{#3}}
\fm@curs{1eso}{E6B84A}{navy}
\fm@curs{2eso}{3B82F6}{white}
\fm@curs{3eso}{FF7A50}{white}
\fm@curs{4eso-a}{A78BFA}{white}
\fm@curs{4eso-b}{6D28D9}{white}
\fm@curs{1batx-mi}{9B1C1C}{white}
\fm@curs{1batx-ccss}{EF6B6B}{white}
\fm@curs{2batx-mii}{1E6B47}{white}
\fm@curs{2batx-ccss}{2BB3A3}{white}

\newcommand{\fmcurs}[1]{%
  \@ifundefined{fm@color@#1}{\PackageError{fentmates}{Curs desconegut: #1}{}}{}%
  \edef\fm@tmp{\noexpand\definecolor{curs}{HTML}{\@nameuse{fm@color@#1}}%
               \noexpand\colorlet{curstext}{\@nameuse{fm@text@#1}}}\fm@tmp
  \colorlet{cursfosc}{curs!72!black}%
  \def\fm@logo{\fm@logodir/vertical-#1.pdf}%
  \def\fm@marca{\fm@logodir/marca-#1.pdf}}

% ---------- capçalera de la primera pàgina ----------
% #1 matèria (petit), #2 curs, #3 tipus de document i unitat, #4 títol
\newcommand{\capcalera}[4]{%
  \gdef\fm@cursnom{#2}%
  \par\noindent
  \begin{tcolorbox}[colback=curs,colframe=curs,boxrule=0pt,arc=2.2mm,width=3.5cm,height=2.15cm,
      valign=center,left=3mm,right=2mm,top=1mm,bottom=1mm,nobeforeafter,box align=center]
    \color{curstext}{\footnotesize #1}\par\vspace{1pt}{\fontsize{19}{21}\selectfont\bfseries #2\par}%
  \end{tcolorbox}\hspace{5mm}%
  \begin{minipage}[c]{\dimexpr\linewidth-3.5cm-5mm-2.9cm\relax}
    {\color{cursfosc}\small\bfseries\MakeUppercase{#3}}\par\vspace{2pt}
    {\color{navy}\fontsize{20}{23}\selectfont\bfseries #4\par}
  \end{minipage}\hfill
  \begin{minipage}[c]{2.6cm}\raggedleft\includegraphics[width=2.6cm]{\fm@logo}\end{minipage}\par
  \vspace{2pt}}

% ---------- barres de secció del color del curs ----------
\renewcommand{\seccio}[2]{\par\needspace{\fmneed\baselineskip}\bigskip
  \begin{tcolorbox}[colback=curs,colframe=curs,coltext=curstext,arc=2mm,boxrule=0pt,left=3mm,right=3mm,top=1.2mm,bottom=1.2mm]
  {\large\bfseries #1\quad #2}
  \end{tcolorbox}\vspace{-2pt}}

% ---------- estrelles de dificultat del color del curs (a Batxillerat no se'n posen) ----------
\@ifundefined{estrelles}{}{%
  \renewcommand{\estrelles}[1]{{\color{curs}\ifcase#1\or\faStar\faStar[regular]\faStar[regular]\or\faStar\faStar\faStar[regular]\or\faStar\faStar\faStar\fi}}}

% ---------- capçalera i peu de les pàgines ----------
% primera pàgina: logos institucionals al peu
\fancypagestyle{primera}{\fancyhf{}\renewcommand{\headrulewidth}{0pt}
  \fancyfoot[L]{\raisebox{-0.32\height}{\includegraphics[height=0.78cm]{logo-gva.png}}\hspace{5mm}%
    \raisebox{-0.32\height}{\includegraphics[height=0.78cm]{logo-vermellar.png}}}
  \fancyfoot[R]{\scriptsize\color{muted}\href{https://creativecommons.org/licenses/by-nc-sa/4.0/deed.ca}{\ccbyncsa}\ \ CC BY-NC-SA 4.0 · fentmates.org\qquad\small\thepage}}
% la resta: el símbol a l'esquerra i el curs a la dreta, amb el títol al centre com abans
\fancyhead[L]{\raisebox{-0.3\height}{\includegraphics[height=0.62cm]{\fm@marca}}}
\fancyhead[R]{\small\bfseries\color{cursfosc}\fm@cursnom}
\renewcommand{\headrule}{\color{curs!45}\hrule height \headrulewidth}
\makeatother

\fmcurs{4eso-b}

\begin{document}
\thispagestyle{primera}

\capcalera{Matemàtiques B}{4t ESO}{Solucions de la fitxa · Unitat 6}{\raggedright Vectors i rectes}
\vspace{8pt}

\begin{tcolorbox}[colback=soft,colframe=soft,arc=2mm,left=3mm,right=3mm,top=1.5mm,bottom=1.5mm]
\small Intenta resoldre cada exercici abans de mirar-ne la solució. Quan una recta es pot escriure de moltes maneres (amb un altre punt o un altre vector director), la teua resposta pot ser diferent i igualment correcta: comprova-ho passant-la a forma general. Pots comprovar molts resultats amb les eines interactives dels apunts de la unitat i amb el laboratori del triangle.
\end{tcolorbox}

% =====================================================
\seccio{1}{Vectors}

\exf{1} \begin{tasks}(2)
  \task $\overrightarrow{AB}=(4-2,\,-1-5)=\res{(2,\,-6)}$
  \task $\overrightarrow{AB}=(1-(-3),\,6-2)=\res{(4,\,4)}$
  \task $\overrightarrow{AB}=(-5-0,\,-1-(-4))=\res{(-5,\,3)}$
  \task $\overrightarrow{AB}=(3-3,\,-2-3)=\res{(0,\,-5)}$
\end{tasks}

\ex{1} \begin{tasks}(2)
  \task $\sqrt{6^2+8^2}=\sqrt{100}=\res{10}$
  \task $\sqrt{(-5)^2+12^2}=\sqrt{169}=\res{13}$
  \task $\sqrt{2^2+(-3)^2}=\sqrt{13}=\res{\sqrt{13}}\approx3,61$
  \task $\sqrt{(-4)^2+(-4)^2}=\sqrt{32}=\res{4\sqrt{2}}\approx5,66$
  \task $\sqrt{0^2+(-7)^2}=\sqrt{49}=\res{7}$
  \task $\sqrt{8^2+(-6)^2}=\sqrt{100}=\res{10}$
\end{tasks}

\ex{1} \begin{enumerate}[label=\alph*),leftmargin=*,itemsep=2pt,topsep=2pt]\item $\overrightarrow{AB}=(-4-3,\,1-(-2))=\res{(-7,\,3)}$; \ $\overrightarrow{BC}=(5-(-4),\,2-1)=\res{(9,\,1)}$\item $|\overrightarrow{AB}|=\sqrt{49+9}=\res{\sqrt{58}}\approx7,62$\item $\overrightarrow{AC}=(2,\,4)$, \ $d(A,C)=\sqrt{4+16}=\sqrt{20}=\res{2\sqrt5}\approx4,47$\end{enumerate}

\ex{1} \begin{enumerate}[label=\alph*),leftmargin=*,itemsep=2pt,topsep=2pt]\item $B=A+\vec v$: \ $\res{(1,\,1)}$; \ $\res{(-1,\,-2)}$; \ $\res{(2,\,-4)}$; \ $\res{(2,\,1)}$\item $A=B-\vec v$: \ $\res{(3,\,-3)}$; \ $\res{(-5,\,3)}$; \ $\res{(-3,\,8)}$; \ $\res{(-1,\,-8)}$\end{enumerate}

\ex{2} $O(0,0)\to\res{(4,-3)}$; \ $(-1,5)\to\res{(3,2)}$; \ $\res{(-4,3)}\to O(0,0)$. Tots són equipol·lents: el mateix vector lliure.

\ex{1} \begin{enumerate}[label=\alph*),leftmargin=*,itemsep=2pt,topsep=2pt]\item \begin{itemize}[leftmargin=*,itemsep=1pt,topsep=2pt]\item $\vec a$: $(-5,\,1)\to(-2,\,3)$, \ $\vec a=\res{(3,\,2)}$, \ $|\vec a|=\sqrt{13}\approx3,61$\item $\vec b$: $(1,\,3)\to(4,\,5)$, \ $\vec b=\res{(3,\,2)}$, \ $|\vec b|=\sqrt{13}\approx3,61$\item $\vec c$: $(-4,\,-1)\to(-4,\,-4)$, \ $\vec c=\res{(0,\,-3)}$, \ $|\vec c|=3$\item $\vec d$: $(5,\,-1)\to(2,\,-3)$, \ $\vec d=\res{(-3,\,-2)}$, \ $|\vec d|=\sqrt{13}\approx3,61$\item $\vec e$: $(-1,\,-3)\to(3,\,-4)$, \ $\vec e=\res{(4,\,-1)}$, \ $|\vec e|=\sqrt{17}\approx4,12$\end{itemize}\item $\vec a$ i $\vec b$ són \res{\text{equipol·lents}} (components $(3,2)$); l'oposat de $\vec a$ és \res{\vec d=(-3,-2)}.\end{enumerate}

\ex{2} 

\begin{center}\begin{tikzpicture}[scale=0.32,baseline=(current bounding box.center)]
\draw[muted!25,very thin] (-4,-4) grid (5,6);
\draw[-latex,muted] (-4.3,0)--(5.5,0) node[right,font=\scriptsize]{$x$}; \draw[-latex,muted] (0,-4.3)--(0,6.5) node[above,font=\scriptsize]{$y$};
\foreach \i in {-4,-2,...,5}{\ifnum\i=0\else\node[font=\tiny,text=muted,below] at (\i,0) {$\i$};\fi}
\foreach \i in {-4,-2,...,6}{\ifnum\i=0\else\node[font=\tiny,text=muted,left] at (0,\i) {$\i$};\fi}
\begin{scope}\clip (-4,-4) rectangle (5,6);
\filldraw[fill=teal!15,draw=teal,thick] (-2,1)--(3,4)--(2,-2)--cycle;
\end{scope}
\draw[-{Latex[length=2.2mm]},accent,line width=1.1pt] (-2,1)--(3,4);
\draw[-{Latex[length=2.2mm]},accent,line width=1.1pt] (3,4)--(2,-2);
\draw[-{Latex[length=2.2mm]},accent,line width=1.1pt] (2,-2)--(-2,1);
\fill[navy] (-2,1) circle (2.4pt) node[left,font=\scriptsize]{$A$};
\fill[navy] (3,4) circle (2.4pt) node[above,font=\scriptsize]{$B$};
\fill[navy] (2,-2) circle (2.4pt) node[below,font=\scriptsize]{$C$};
\end{tikzpicture}\end{center}
$\overrightarrow{AB}=(3-(-2),\,4-1)=\res{(5,\,3)}$; \ $\overrightarrow{BC}=(2-3,\,-2-4)=\res{(-1,\,-6)}$; \ $\overrightarrow{CA}=(-2-2,\,1-(-2))=\res{(-4,\,3)}$. La suma és $\res{\vec 0=(0,0)}$: si anem de $A$ a $B$, de $B$ a $C$ i de $C$ a $A$, tornem al punt de partida.


% =====================================================
\seccio{2}{Operacions amb vectors}

\exf{1} \begin{tasks}(3)
  \task $\vec u+\vec v=\res{(1,\,1)}$
  \task $\vec u-\vec v=\res{(3,\,-7)}$
  \task $3\vec u=\res{(6,\,-9)}$
  \task $-2\vec v=\res{(2,\,-8)}$
  \task $2\vec u+3\vec v=\res{(1,\,6)}$
  \task $\tfrac12\vec u-\vec v=\res{(2,\,-\frac{11}{2})}$
\end{tasks}

\ex{1} \begin{tasks}(2)
  \task $\vec a+\vec b+\vec c=\res{(4,\,-1)}$
  \task $2\vec a-\vec b+\vec c=\res{(9,\,-11)}$
  \task $-\vec a+3\vec b-2\vec c=\res{(-17,\,13)}$
  \task $\tfrac15\vec c-\tfrac12\vec b=\res{(2,\,-\frac{3}{2})}$
\end{tasks}

\ex{2} 

\begin{center}\begin{tikzpicture}[scale=0.55,baseline=(current bounding box.center)]
\draw[muted!25,very thin] (-2,-3) grid (7,4);
\draw[-latex,muted] (-2.3,0)--(7.5,0) node[right,font=\scriptsize]{$x$}; \draw[-latex,muted] (0,-3.3)--(0,4.5) node[above,font=\scriptsize]{$y$};
\foreach \i in {-2,0,...,7}{\ifnum\i=0\else\node[font=\tiny,text=muted,below] at (\i,0) {$\i$};\fi}
\foreach \i in {-2,0,...,4}{\ifnum\i=0\else\node[font=\tiny,text=muted,left] at (0,\i) {$\i$};\fi}
\begin{scope}\clip (-2,-3) rectangle (7,4);
\end{scope}
\draw[-{Latex[length=2.2mm]},accent,line width=1.1pt] (0,0)--(3,1);
\node[accent,font=\small,below right] at (1.5,0.5) {$\vec u$};
\draw[-{Latex[length=2.2mm]},teal,line width=1.1pt] (0,0)--(-1,2);
\node[teal,font=\small,left] at (-0.5,1) {$\vec v$};
\draw[-{Latex[length=2.2mm]},teal,line width=1.1pt] (3,1)--(2,3);
\draw[-{Latex[length=2.2mm]},accent,line width=1.1pt] (-1,2)--(2,3);
\draw[-{Latex[length=2.2mm]},navy,line width=1.1pt] (0,0)--(2,3);
\node[navy,font=\small,left] at (1,1.5) {$\vec u+\vec v$};
\draw[-{Latex[length=2.2mm]},violet,line width=1.1pt] (-1,2)--(3,1);
\node[violet,font=\small,above right] at (1,1.5) {$\vec u-\vec v$};
\draw[-{Latex[length=2.2mm]},orange,line width=1.1pt] (0,-2)--(6,0);
\node[orange,font=\small,below right] at (3,-1) {$2\vec u$};
\end{tikzpicture}\end{center}
$\vec u+\vec v=\res{(2,3)}$ (diagonal del paral·lelogram); \ $\vec u-\vec v=\res{(4,-1)}$ (va de la punta de $\vec v$ a la de $\vec u$); \ $2\vec u=\res{(6,2)}$ (en taronja, dibuixat a partir de $(0,-2)$).

\ex{1} \begin{tasks}(2)
  \task $\dfrac{2}{-1}=\dfrac{-6}{3}$: $\res{\text{sí}}$ \ ($\vec v=-\frac{1}{2}\vec u$)
  \task $\dfrac{4}{6}=\dfrac{6}{9}$: $\res{\text{sí}}$ \ ($\vec v=\frac{3}{2}\vec u$)
  \task $\dfrac{3}{5}\neq\dfrac{5}{3}$: $\res{\text{no}}$
  \task $\dfrac{-2}{1}\neq\dfrac{8}{4}$: $\res{\text{no}}$
  \task $0\cdot(-2)=5\cdot0$: $\res{\text{sí}}$ \ ($\vec v=-\frac{2}{5}\vec u$)
  \task $\dfrac{10}{-15}=\dfrac{-4}{6}$: $\res{\text{sí}}$ \ ($\vec v=-\frac{3}{2}\vec u$)
\end{tasks}

\ex{2} \begin{tasks}(2)
  \task $\dfrac{x}{12}=\dfrac{3}{x}\Rightarrow x^2=36\Rightarrow \res{x=\pm6}$
  \task $\dfrac{2}{x}=\dfrac{-5}{10}\Rightarrow \res{x=-4}$
  \task $4x=3(x+1)\Rightarrow \res{x=3}$
\end{tasks}

\ex{2} $\vec u=-\tfrac13\vec v=\res{(1,-3)}$; \ $\res{\vec v=-3\vec u}$.

\ex{3} $3\vec x=2\vec v-\vec u=(4,-2)-(1,-5)=(3,3)\Rightarrow\vec x=\res{(1,\,1)}$


% =====================================================
\seccio{3}{Producte escalar. Vectors perpendiculars}

\exf{1} \begin{tasks}(2)
  \task $3\cdot1+2\cdot(-4)=\res{-5}$ \ (angle obtús)
  \task $-2\cdot4+5\cdot(-1)=\res{-13}$ \ (angle obtús)
  \task $6\cdot1+(-3)\cdot2=\res{0}$ \ (perpendiculars)
  \task $0\cdot(-5)+3\cdot7=\res{21}$ \ (angle agut)
\end{tasks}

\ex{1} $\overrightarrow{AB}=(-7,9)$, \ $\overrightarrow{CD}=(-4,-6)$, \ $\overrightarrow{DA}=(5,4)$.\begin{tasks}(2)
  \task $\res{(-11,\,3)}$
  \task $(-7,9)-(-12,-18)=\res{(5,\,27)}$
  \task $(-14,18)\cdot(-4,-6)=56-108=\res{-52}$
  \task $(-35,45)\cdot(-5,-4)=175-180=\res{-5}$
  \task $\sqrt{49+81}=\res{\sqrt{130}}\approx11,4$
  \task $\sqrt{16+36}=\sqrt{52}=\res{2\sqrt{13}}\approx7,21$
\end{tasks}

\ex{1} \begin{tasks}(2)
  \task $\vec u\cdot\vec v=2\cdot(-1)+1\cdot2=0$: $\res{\text{perpendiculars}}$
  \task $-6\cdot(-2)=4\cdot3$: $\res{\text{paral·lels}}$
  \task $\vec u\cdot\vec v=20\neq0$ i no proporcionals: $\res{\text{cap}}$
  \task $\vec u\cdot\vec v=4\cdot9+(-6)\cdot6=0$: $\res{\text{perpendiculars}}$
  \task $1\cdot(-3)=1\cdot(-3)$: $\res{\text{paral·lels}}$
  \task $\vec u\cdot\vec v=7\cdot0+0\cdot(-2)=0$: $\res{\text{perpendiculars}}$
\end{tasks}

\ex{2} Perpendiculars: $\res{(2,5)}$ i $\res{(-2,-5)}$ ($10-10=0$). Paral·lels: $\res{(10,-4)}$ (mateix sentit) i $\res{(-5,2)}$ (sentit contrari). \ $\vec v\cdot\vec w=v_1(-v_2)+v_2v_1=\res{0}$.

\ex{2} \begin{tasks}(2)
  \task $3(x+1)-2x=0\Rightarrow \res{x=-3}$
  \task $6x-12=0\Rightarrow \res{x=2}$
  \task $x^2-16=0\Rightarrow \res{x=\pm4}$
\end{tasks}

\ex{2} 

\begin{center}\begin{tikzpicture}[scale=0.42,baseline=(current bounding box.center)]
\draw[muted!25,very thin] (-5,-2) grid (3,3);
\draw[-latex,muted] (-5.3,0)--(3.5,0) node[right,font=\scriptsize]{$x$}; \draw[-latex,muted] (0,-2.3)--(0,3.5) node[above,font=\scriptsize]{$y$};
\foreach \i in {-4,-2,...,3}{\ifnum\i=0\else\node[font=\tiny,text=muted,below] at (\i,0) {$\i$};\fi}
\foreach \i in {-2,0,...,3}{\ifnum\i=0\else\node[font=\tiny,text=muted,left] at (0,\i) {$\i$};\fi}
\begin{scope}\clip (-5,-2) rectangle (3,3);
\end{scope}
\draw[-{Latex[length=2.2mm]},accent,line width=1.1pt] (0,0)--(2,-1);
\node[accent,font=\small,below] at (1,-0.5) {$\vec a$};
\draw[-{Latex[length=2.2mm]},teal,line width=1.1pt] (0,0)--(1,2);
\node[teal,font=\small,left] at (0.5,1) {$\vec b$};
\draw[-{Latex[length=2.2mm]},navy,line width=1.1pt] (0,0)--(-4,2);
\node[navy,font=\small,above] at (-2,1) {$\vec c$};
\end{tikzpicture}\end{center}
$\vec a\cdot\vec b=2-2=0$: $\res{\vec a\perp\vec b}$; \ $\vec c=-2\vec a$: $\res{\vec a\parallel\vec c}$ (sentit contrari); \ $\vec b\cdot\vec c=-4+4=0$: $\res{\vec b\perp\vec c}$.

\ex{3} Un perpendicular és $(-4,3)$, de mòdul 5. Multiplicant per $\pm2$: $\res{(-8,6)}$ i $\res{(8,-6)}$.


% =====================================================
\seccio{4}{Aplicacions dels vectors}

\exf{1} 

\begin{center}\begin{tikzpicture}[scale=0.42,baseline=(current bounding box.center)]
\draw[muted!25,very thin] (-3,-2) grid (6,3);
\draw[-latex,muted] (-3.3,0)--(6.5,0) node[right,font=\scriptsize]{$x$}; \draw[-latex,muted] (0,-2.3)--(0,3.5) node[above,font=\scriptsize]{$y$};
\foreach \i in {-2,0,...,6}{\ifnum\i=0\else\node[font=\tiny,text=muted,below] at (\i,0) {$\i$};\fi}
\foreach \i in {-2,0,...,3}{\ifnum\i=0\else\node[font=\tiny,text=muted,left] at (0,\i) {$\i$};\fi}
\begin{scope}\clip (-3,-2) rectangle (6,3);
\end{scope}
\draw[-{Latex[length=2.2mm]},accent,line width=1.1pt] (4,0)--(-1,1);
\node[accent,font=\small,above] at (1.5,0.5) {$\overrightarrow{AB}$};
\fill[navy] (4,0) circle (2.4pt) node[below,font=\scriptsize]{$A$};
\fill[navy] (-1,1) circle (2.4pt) node[above left,font=\scriptsize]{$B$};
\end{tikzpicture}\end{center}
$\overrightarrow{AB}=(-5,1)$, \ $d(A,B)=\sqrt{25+1}=\res{\sqrt{26}}\approx5,1$; \ perpendicular: $\res{(1,5)}$; \ paral·lel: $\res{(-10,2)}$.

\ex{1} \begin{tasks}(2)
  \task $\sqrt{3^2+4^2}=\sqrt{25}=\res{5}$
  \task $\sqrt{5^2+(-12)^2}=\sqrt{169}=\res{13}$
  \task $\sqrt{(-4)^2+0^2}=\sqrt{16}=\res{4}$
  \task $\sqrt{3^2+3^2}=\sqrt{18}=\res{3\sqrt{2}}\approx4,24$
  \task $\sqrt{(-5)^2+5^2}=\sqrt{50}=\res{5\sqrt{2}}\approx7,07$
  \task $\sqrt{6^2+8^2}=\sqrt{100}=\res{10}$
\end{tasks}

\ex{1} \begin{tasks}(3)
  \task $\left(\frac{2+8}{2},\frac{6+0}{2}\right)=\res{(5,\,3)}$
  \task $\left(\frac{-3+5}{2},\frac{4+(-2)}{2}\right)=\res{(1,\,1)}$
  \task $\left(\frac{1+(-2)}{2},\frac{-5+2}{2}\right)=\res{(-\frac{1}{2},\,-\frac{3}{2})}$
  \task $\left(\frac{0+(-6)}{2},\frac{7+(-3)}{2}\right)=\res{(-3,\,2)}$
  \task $\left(\frac{2+(-4)}{2},\frac{1+3}{2}\right)=\res{(-1,\,2)}$
\end{tasks}

\ex{1} \begin{tasks}(2)
  \task $B=A+2\overrightarrow{AM}=(1,\,4)+2(2,\,-3)=\res{(5,\,-2)}$
  \task $B=A+2\overrightarrow{AM}=(-2,\,0)+2(2,\,3)=\res{(2,\,6)}$
  \task $B=A+2\overrightarrow{AM}=(4,\,-3)+2(-3,\,2)=\res{(-2,\,1)}$
  \task $B=A+2\overrightarrow{AM}=(-5,\,2)+2(4,\,-4)=\res{(3,\,-6)}$
\end{tasks}

\ex{1} $B$ és el punt mig de $AA'$: $\frac{-1+x}{2}=2$, $\frac{0+y}{2}=-8$ $\Rightarrow A'\res{(5,-16)}$. \ $P'=P+2\overrightarrow{PQ}=(3,-2)+2(-2,3)=\res{(-1,4)}$.

\ex{2} \begin{enumerate}[label=\alph*),leftmargin=*,itemsep=2pt,topsep=2pt]\item $\overrightarrow{AB}=(2,1)$, $\overrightarrow{AC}=(-2,-1)$: $\frac{2}{-2}=\frac{1}{-1}$: $\res{\text{sí que estan alineats}}$\item $\overrightarrow{AB}=(-1,2)$, $\overrightarrow{AC}=(1,k-1)$: $\frac{1}{-1}=\frac{k-1}{2}\Rightarrow k-1=-2\Rightarrow \res{k=-1}$\item $\overrightarrow{PQ}=(2,3)$, $\overrightarrow{QR}=(3,2)$: $\frac23\neq\frac32$: $\res{\text{no estan alineats}}$\end{enumerate}

\ex{2} \begin{enumerate}[label=\alph*),leftmargin=*,itemsep=2pt,topsep=2pt]\item $\overrightarrow{AB}=\overrightarrow{DC}=(3,2)\Rightarrow B=A+\overrightarrow{DC}=\res{(5,3)}$\item $\overrightarrow{AD}=\overrightarrow{BC}=(3,-3)\Rightarrow D=\res{(1,-2)}$\end{enumerate}

\ex{2} $\overrightarrow{AB}=(9,-3)$, $\overrightarrow{AC}=(2,6)$: $\overrightarrow{AB}\cdot\overrightarrow{AC}=18-18=0$. És rectangle en $\res{A}$.

\ex{2} $\overrightarrow{AB}=(1,4)$ i $\overrightarrow{BC}=(4,-1)$ són perpendiculars ($4-4=0$) i tenen el mateix mòdul ($\sqrt{17}$): $A$, $B$ i $C$ són consecutius. \ $D=A+\overrightarrow{BC}=\res{(8,-3)}$.

\ex{2} $\overrightarrow{AB}=(3,1)$, $\overrightarrow{BC}=(-1,3)$, $\overrightarrow{CD}=(-3,-1)$, $\overrightarrow{DA}=(1,-3)$: tots tenen mòdul $\sqrt{10}$ i $\overrightarrow{AB}\cdot\overrightarrow{BC}=-3+3=0$: és un quadrat. $P=\res{4\sqrt{10}\approx12,65\ \text{u}}$; \ $S=(\sqrt{10})^2=\res{10\ \text{u}^2}$.

\ex{2} Costats: $(-2,-3)$, $(2,-3)$, $(2,3)$, $(-2,3)$, tots de mòdul $\sqrt{13}$. \ Diagonals: $|\overrightarrow{AC}|=|(0,-6)|=\res{6}$, \ $|\overrightarrow{BD}|=|(4,0)|=\res{4}$. \ $S=\frac{6\cdot4}{2}=\res{12\ \text{u}^2}$.

\ex{2} $\overrightarrow{AB}=(6,3)$, $|\overrightarrow{AB}|=\sqrt{45}=3\sqrt5$; \ $\overrightarrow{BC}=(-3,3)$, $|\overrightarrow{BC}|=\sqrt{18}=3\sqrt2$; \ $\overrightarrow{CA}=(-3,-6)$, $|\overrightarrow{CA}|=3\sqrt5$. \ $P=\res{6\sqrt5+3\sqrt2\approx17,66\ \text{u}}$; \ és $\res{\text{isòsceles}}$.

\ex{2} \begin{enumerate}[label=\alph*),leftmargin=*,itemsep=2pt,topsep=2pt]\item El centre és el punt mig: $C=\res{(1,1)}$; \ $r=d(C,Q)=\sqrt{16+9}=\res{5}$.\item Com que els radis són iguals, és el punt mig de $AB$: $\res{(2,1)}$.\end{enumerate}

\ex{3} $P(x,0)$: $(x-1)^2+9=(x-5)^2+1\Rightarrow -2x+10=-10x+26\Rightarrow 8x=16\Rightarrow x=2$. \ $\res{P(2,0)}$ \ ($d=\sqrt{10}$).

\ex{3} $M(1,3)$, $N(4,4)$, $\overrightarrow{MN}=(3,1)=\tfrac12(6,2)=\res{\tfrac12\overrightarrow{AB}}$: paral·lels i de mòdul la meitat ($\sqrt{10}$ i $2\sqrt{10}$).


% =====================================================
\seccio{5}{Equacions vectorial, paramètriques i contínua}

\exf{1} \begin{tasks}[after-item-skip=6pt](2)
  \task $(x,y)=(2,\,-1)+\lambda(3,\,4)$\\ $\left\{\begin{array}{l}x=2+3\lambda\\ y=-1+4\lambda\end{array}\right.$\quad $\res{\dfrac{x-2}{3}=\dfrac{y+1}{4}}$
  \task $(x,y)=(0,\,5)+\lambda(-2,\,1)$\\ $\left\{\begin{array}{l}x=-2\lambda\\ y=5+\lambda\end{array}\right.$\quad $\res{\dfrac{x}{-2}=\dfrac{y-5}{1}}$
  \task $(x,y)=(-4,\,-3)+\lambda(1,\,-5)$\\ $\left\{\begin{array}{l}x=-4+\lambda\\ y=-3-5\lambda\end{array}\right.$\quad $\res{\dfrac{x+4}{1}=\dfrac{y+3}{-5}}$
  \task $(x,y)=(6,\,0)+\lambda(-3,\,-3)$\\ $\left\{\begin{array}{l}x=6-3\lambda\\ y=-3\lambda\end{array}\right.$\quad $\res{\dfrac{x-6}{-3}=\dfrac{y}{-3}}$
\end{tasks}

\ex{1} \begin{tasks}[after-item-skip=6pt](2)
  \task $\overrightarrow{AB}=(3,\,-3)\parallel(1,\,-1)$\\ $(x,y)=(1,\,2)+\lambda(1,\,-1)$\\ $\left\{\begin{array}{l}x=1+\lambda\\ y=2-\lambda\end{array}\right.$\quad $\res{\dfrac{x-1}{1}=\dfrac{y-2}{-1}}$
  \task $\overrightarrow{AB}=(5,\,4)$\\ $(x,y)=(-2,\,0)+\lambda(5,\,4)$\\ $\left\{\begin{array}{l}x=-2+5\lambda\\ y=4\lambda\end{array}\right.$\quad $\res{\dfrac{x+2}{5}=\dfrac{y}{4}}$
  \task $\overrightarrow{AB}=(-6,\,-3)\parallel(2,\,1)$\\ $(x,y)=(5,\,-3)+\lambda(2,\,1)$\\ $\left\{\begin{array}{l}x=5+2\lambda\\ y=-3+\lambda\end{array}\right.$\quad $\res{\dfrac{x-5}{2}=\dfrac{y+3}{1}}$
  \task $\overrightarrow{AB}=(-4,\,6)\parallel(2,\,-3)$\\ $(x,y)=(0,\,0)+\lambda(2,\,-3)$\\ $\left\{\begin{array}{l}x=2\lambda\\ y=-3\lambda\end{array}\right.$\quad $\res{\dfrac{x}{2}=\dfrac{y}{-3}}$
\end{tasks}

\ex{1} (Hi ha moltes respostes correctes: depenen del punt i del vector que es prenga.)\begin{itemize}[leftmargin=*,itemsep=1pt,topsep=2pt]\item $r$: passa per $(-4,-3)$ i $(2,1)$: \ $\res{(x,y)=(-1,\,-1)+\lambda(3,\,2)}$, \ $\left\{\begin{array}{l}x=-1+3\lambda\\ y=-1+2\lambda\end{array}\right.$\item $s$: passa per $(0,4)$ i $(1,2)$: \ $\res{(x,y)=(0,\,4)+\lambda(1,\,-2)}$, \ $\left\{\begin{array}{l}x=\lambda\\ y=4-2\lambda\end{array}\right.$\item $t$: passa per $(-5,5)$ i $(1,3)$: \ $\res{(x,y)=(1,\,3)+\lambda(3,\,-1)}$, \ $\left\{\begin{array}{l}x=1+3\lambda\\ y=3-\lambda\end{array}\right.$\item $u$: vertical, passa per $(4,0)$: \ $\res{(x,y)=(4,\,0)+\lambda(0,\,1)}$, \ $\left\{\begin{array}{l}x=4\\ y=\lambda\end{array}\right.$. El vector director $(0,1)$ té una component nul·la i no podem dividir per 0: la recta és $x=4$.\end{itemize}

\ex{2} \begin{tasks}[after-item-skip=4pt](2)
  \task $3=1+\lambda\Rightarrow\lambda=2$; \ $-1=3-2\lambda\Rightarrow\lambda=2$: el mateix $\lambda$, $\res{P\in r}$
  \task $3=5-2\lambda\Rightarrow\lambda=1$; \ $-1=2+\lambda\Rightarrow\lambda=-3$: $\res{P\notin r}$
  \task $\frac{3-7}{4}=-1$, \ $\frac{-1+3}{-2}=-1$: $\res{P\in r}$
  \task $\frac{3+1}{2}=2$, \ $\frac{-1-1}{3}=-\frac23$: $\res{P\notin r}$
\end{tasks}

\ex{2} \begin{enumerate}[label=\alph*),leftmargin=*,itemsep=2pt,topsep=2pt]\item $\lambda=0,1,-1$: $\res{(3,1)}$, $\res{(5,-3)}$, $\res{(1,5)}$\item $-7=1-4\lambda\Rightarrow\lambda=2\Rightarrow k=3+4=\res{7}$\item $\res{\dfrac{x-3}{2}=\dfrac{y-1}{-4}}$; \ $\frac{7-3}{2}=2=\frac{-7-1}{-4}$\checkmark\end{enumerate}

\ex{2} Paral·lela a $OX$: $(x,y)=(-2,5)+\lambda(1,0)$, és a dir, $\res{y=5}$. \ Paral·lela a $OY$: $(x,y)=(-2,5)+\lambda(0,1)$, és a dir, $\res{x=-2}$. \ Eixos: $\res{y=0}$ (eix $OX$) i $\res{x=0}$ (eix $OY$).


% =====================================================
\seccio{6}{Pendent. Equacions punt-pendent, explícita i general}

\exf{1} \begin{tasks}[after-item-skip=4pt](2)
  \task $m=\dfrac{8-2}{3-1}=\res{3}$ (creixent)
  \task $m=\dfrac{2-4}{2-(-2)}=\res{-\frac{1}{2}}$ (decreixent)
  \task $m=\dfrac{-3-(-3)}{5-0}=\res{0}$ (horitzontal)
  \task $\Delta x=0$: recta vertical, $\res{\text{no té pendent}}$
\end{tasks}

\ex{1} \begin{tasks}[after-item-skip=4pt](2)
  \task $y+2=4(x-1)$; \ $\res{y=4x-6}$
  \task $y-5=-3x$; \ $\res{y=-3x+5}$
  \task $y-1=\frac{1}{2}(x+4)$; \ $\res{y=\frac{1}{2}x+3}$
  \task $y=-\frac{2}{3}(x-3)$; \ $\res{y=-\frac{2}{3}x+2}$
\end{tasks}

\ex{1} \begin{enumerate}[label=\alph*),leftmargin=*,itemsep=2pt,topsep=2pt]\item $\res{m=-1}$\item $\res{(1,-1)}$\item $\res{(0,2)}$ i $\res{(2,0)}$

\begin{center}\begin{tikzpicture}[scale=0.3,baseline=(current bounding box.center)]
\draw[muted!25,very thin] (-3,-2) grid (4,4);
\draw[-latex,muted] (-3.3,0)--(4.5,0) node[right,font=\scriptsize]{$x$}; \draw[-latex,muted] (0,-2.3)--(0,4.5) node[above,font=\scriptsize]{$y$};
\foreach \i in {-2,0,...,4}{\ifnum\i=0\else\node[font=\tiny,text=muted,below] at (\i,0) {$\i$};\fi}
\foreach \i in {-2,0,...,4}{\ifnum\i=0\else\node[font=\tiny,text=muted,left] at (0,\i) {$\i$};\fi}
\begin{scope}\clip (-3,-2) rectangle (4,4);
\draw[accent,thick] (-3,5)--(4,-2);
\end{scope}
\fill[navy] (0,2) circle (2.4pt) node[right,font=\scriptsize]{$(0,2)$};
\fill[navy] (2,0) circle (2.4pt) node[above right,font=\scriptsize]{$(2,0)$};
\end{tikzpicture}\end{center}
\item $-(-1)+2=3\neq4$: $\res{A\notin r}$\end{enumerate}

\ex{1} \begin{multicols}{2}\begin{itemize}[leftmargin=*,itemsep=1pt,topsep=2pt]\item Vectorial: $(x,y)=(2,\,2)+\lambda(-4,\,2)$\item Paramètriques: $\left\{\begin{array}{l}x=2-4\lambda\\ y=2+2\lambda\end{array}\right.$\item Contínua: $\dfrac{x-2}{-4}=\dfrac{y-2}{2}$\item Punt-pendent: $y-2=-\frac{1}{2}(x-2)$\item Explícita: $y=-\frac{1}{2}x+3$\item General: $\res{x+2y-6=0}$\end{itemize}\end{multicols}

\ex{1} $\overrightarrow{PQ}=(2,2)\parallel(1,1)$, $m=1$.\begin{multicols}{2}\begin{itemize}[leftmargin=*,itemsep=1pt,topsep=2pt]\item Vectorial: $(x,y)=(-3,\,-4)+\lambda(1,\,1)$\item Paramètriques: $\left\{\begin{array}{l}x=-3+\lambda\\ y=-4+\lambda\end{array}\right.$\item Contínua: $\dfrac{x+3}{1}=\dfrac{y+4}{1}$\item Punt-pendent: $y+4=(x+3)$\item Explícita: $y=x-1$\item General: $\res{x-y-1=0}$\end{itemize}\end{multicols}

\ex{1} Vector director $(1,2)$.\begin{multicols}{2}\begin{itemize}[leftmargin=*,itemsep=1pt,topsep=2pt]\item Vectorial: $(x,y)=(5,\,-1)+\lambda(1,\,2)$\item Paramètriques: $\left\{\begin{array}{l}x=5+\lambda\\ y=-1+2\lambda\end{array}\right.$\item Contínua: $\dfrac{x-5}{1}=\dfrac{y+1}{2}$\item Punt-pendent: $y+1=2(x-5)$\item Explícita: $y=2x-11$\item General: $\res{2x-y-11=0}$\end{itemize}\end{multicols}

\ex{1} $\overrightarrow{AB}=(2,6)$.\begin{multicols}{2}\begin{itemize}[leftmargin=*,itemsep=1pt,topsep=2pt]\item Vectorial: $(x,y)=(1,\,0)+\lambda(2,\,6)$\item Paramètriques: $\left\{\begin{array}{l}x=1+2\lambda\\ y=6\lambda\end{array}\right.$\item Contínua: $\dfrac{x-1}{2}=\dfrac{y}{6}$\item Punt-pendent: $y=3(x-1)$\item Explícita: $y=3x-3$\item General: $\res{3x-y-3=0}$\end{itemize}\end{multicols}

\ex{2} \begin{enumerate}[label=\alph*),leftmargin=*,itemsep=2pt,topsep=2pt]\item $\left\{\begin{array}{l}x=2-t\\ y=-3+2t\end{array}\right.$:\begin{multicols}{2}\begin{itemize}[leftmargin=*,itemsep=1pt,topsep=2pt]\item Vectorial: $(x,y)=(2,\,-3)+t(-1,\,2)$\item Paramètriques: $\left\{\begin{array}{l}x=2-t\\ y=-3+2t\end{array}\right.$\item Contínua: $\dfrac{x-2}{-1}=\dfrac{y+3}{2}$\item Punt-pendent: $y+3=-2(x-2)$\item Explícita: $y=-2x+1$\item General: $\res{2x+y-1=0}$\end{itemize}\end{multicols}\item $(x,y)=(0,3)+t(2,1)$:\begin{multicols}{2}\begin{itemize}[leftmargin=*,itemsep=1pt,topsep=2pt]\item Vectorial: $(x,y)=(0,\,3)+t(2,\,1)$\item Paramètriques: $\left\{\begin{array}{l}x=2t\\ y=3+t\end{array}\right.$\item Contínua: $\dfrac{x}{2}=\dfrac{y-3}{1}$\item Punt-pendent: $y-3=\frac{1}{2}x$\item Explícita: $y=\frac{1}{2}x+3$\item General: $\res{x-2y+6=0}$\end{itemize}\end{multicols}\item $y=3x-1$:\begin{multicols}{2}\begin{itemize}[leftmargin=*,itemsep=1pt,topsep=2pt]\item Vectorial: $(x,y)=(0,\,-1)+t(1,\,3)$\item Paramètriques: $\left\{\begin{array}{l}x=t\\ y=-1+3t\end{array}\right.$\item Contínua: $\dfrac{x}{1}=\dfrac{y+1}{3}$\item Punt-pendent: $y+1=3x$\item Explícita: $y=3x-1$\item General: $\res{3x-y-1=0}$\end{itemize}\end{multicols}\item $y-3=3(x-5)$:\begin{multicols}{2}\begin{itemize}[leftmargin=*,itemsep=1pt,topsep=2pt]\item Vectorial: $(x,y)=(5,\,3)+t(1,\,3)$\item Paramètriques: $\left\{\begin{array}{l}x=5+t\\ y=3+3t\end{array}\right.$\item Contínua: $\dfrac{x-5}{1}=\dfrac{y-3}{3}$\item Punt-pendent: $y-3=3(x-5)$\item Explícita: $y=3x-12$\item General: $\res{3x-y-12=0}$\end{itemize}\end{multicols}\item $2x+y-5=0$:\begin{multicols}{2}\begin{itemize}[leftmargin=*,itemsep=1pt,topsep=2pt]\item Vectorial: $(x,y)=(0,\,5)+t(1,\,-2)$\item Paramètriques: $\left\{\begin{array}{l}x=t\\ y=5-2t\end{array}\right.$\item Contínua: $\dfrac{x}{1}=\dfrac{y-5}{-2}$\item Punt-pendent: $y-5=-2x$\item Explícita: $y=-2x+5$\item General: $\res{2x+y-5=0}$\end{itemize}\end{multicols}\end{enumerate}

\ex{2} \begin{enumerate}[label=\alph*),leftmargin=*,itemsep=2pt,topsep=2pt]\item $-2=3+4t\Rightarrow t=-\frac54$; \ $1=2-3t\Rightarrow t=\frac13$: $\res{M\notin r}$\item $\frac{-2-5}{2}=-\frac72$; \ $\frac{3-1}{4}=\frac12$: $\res{M\notin r}$\item $-6+5+30=29\neq0$: $\res{M\notin r}$\item $\frac{-2}{2}+2=1$\checkmark: $\res{M\in r}$\item $1+2=3$; \ $3(-2+5)=9$: $\res{M\notin r}$\end{enumerate}

\ex{1} \begin{tasks}(2)
  \task $\res{(2,4)}$, \ $\res{(5,-3)}$
  \task $\res{(3,-2)}$, \ $\res{(3,2)}$
  \task $\res{(-1,-8)}$, \ $\res{(9,-6)}$
  \task $\res{(-4,0)}$, \ $\res{(3,2)}$ \ ($\vec v=(-B,A)$)
\end{tasks}

\ex{1} \begin{tasks}(1)
  \task $m=\res{\frac{3}{4}}$, \ $n=\res{3}$, \ $\vec v=\res{(4,\,3)}$, \ $\vec n=\res{(3,\,-4)}$
  \task $m=\res{-2}$, \ $n=\res{5}$, \ $\vec v=\res{(1,\,-2)}$, \ $\vec n=\res{(2,\,1)}$
  \task Vertical ($x=3$): no té pendent ni ordenada en l'origen; \ $\vec v=\res{(0,1)}$, \ $\vec n=\res{(1,0)}$
  \task $m=\res{0}$, \ $n=\res{-2}$, \ $\vec v=\res{(1,\,0)}$, \ $\vec n=\res{(0,\,1)}$
\end{tasks}

\ex{2} \begin{tasks}(1)
  \task $\overrightarrow{AB}=(3,4)$, $\vec n=(4,-3)$: $4x-3y+C=0$, $8+9+C=0\Rightarrow C=-17$: $\res{4x-3y-17=0}$
  \task $\overrightarrow{AB}=(4,0)$: recta horitzontal $\res{y-4=0}$
  \task $\vec n=(1,3)$: $x+3y+C=0$, $-6+C=0\Rightarrow C=6$: $\res{x+3y+6=0}$
  \task $2x-5y+C=0$, $8-5+C=0\Rightarrow C=-3$: $\res{2x-5y-3=0}$
\end{tasks}

\ex{1} \begin{tasks}[after-item-skip=4pt](2)
  \task $m=\frac{1}{2}$: \ $y+1=\frac{1}{2}(x-2)\Rightarrow \res{y=\frac{1}{2}x-2}$
  \task $m=-3$: \ $y-5=-3(x+3)\Rightarrow \res{y=-3x-4}$
  \task $m=-\frac{2}{3}$: \ $y-4=-\frac{2}{3}x\Rightarrow \res{y=-\frac{2}{3}x+4}$
  \task $m=0$: \ $y+2=0\Rightarrow \res{y=-2}$
\end{tasks}

\ex{2} \begin{enumerate}[label=\alph*),leftmargin=*,itemsep=2pt,topsep=2pt]\item $y-0=-\frac32(x-4)\Rightarrow \res{y=-\frac{3}{2}x+6}$\item Pendent 0: $\res{y=-2}$\item $m=-\frac52$, $n=0$: $\res{y=-\frac{5}{2}x}$\end{enumerate}

\ex{3} $y=0\Rightarrow x=4$: $\res{P(4,0)}$; \ $x=0\Rightarrow y=3$: $\res{Q(0,3)}$. \ Triangle rectangle de catets 4 i 3: $S=\frac{4\cdot3}{2}=\res{6\ \text{u}^2}$, \ $d(P,Q)=5$, \ $\res{P=12\ \text{u}}$.


% =====================================================
\seccio{7}{Posició relativa de dues rectes}

\exf{1} \begin{tasks}[after-item-skip=4pt](2)
  \task Sumant: $3x=6\Rightarrow x=2,\ y=3$: \ $\res{(2,3)}$
  \task $x=-8-4y\Rightarrow 3(-8-4y)-2y=4\Rightarrow y=-2,\ x=0$: \ $\res{(0,-2)}$
  \task $x=y+5\Rightarrow 2y+10+3y=0\Rightarrow y=-2,\ x=3$: \ $\res{(3,-2)}$
  \task Sumant: $6x=0\Rightarrow x=0,\ y=\frac32$: \ $\res{\left(0,\tfrac32\right)}$
\end{tasks}

\ex{1} \begin{tasks}[after-item-skip=4pt](2)
  \task $\frac{2}{-4}=\frac{-1}{2}=\frac{3}{-6}$: $\res{\text{coincidents}}$
  \task $\frac12=\frac36\neq\frac{-2}{1}$: $\res{\text{paral·leles}}$
  \task $\frac31\neq\frac{-1}{2}$: $\res{\text{secants}}$ en $\left(\frac37,\frac{16}{7}\right)$
  \task $\frac{1/2}{1}=\frac12=\frac{-1}{-2}$: $\res{\text{coincidents}}$
\end{tasks}

\ex{1} \begin{enumerate}[label=\alph*),leftmargin=*,itemsep=2pt,topsep=2pt]\item $\frac22\neq\frac{3}{-3}$: secants. Sumant: $4x=16\Rightarrow x=4$, $y=-1$: $\res{(4,-1)}$\item $\frac{2}{-2}=\frac{-1}{1}\neq\frac{-3}{-1}$: $\res{\text{paral·leles}}$\item $r:2x-y-1=0$ ($\vec v=(2,4)$), \ $s:x-5y+33=0$ ($\vec v=(5,1)$): secants. \ $y=2x-1\Rightarrow x-10x+5+33=0\Rightarrow x=\frac{38}{9}$, $y=\frac{67}{9}$: $\res{\left(\tfrac{38}{9},\tfrac{67}{9}\right)}$\item $r:3x-2y+2=0$ ($m=\frac32$, $y=\frac32x+1$) i $s:3x-2y-4=0$: mateix pendent i diferent ordenada en l'origen: $\res{\text{paral·leles}}$\end{enumerate}

\ex{2} \begin{enumerate}[label=\alph*),leftmargin=*,itemsep=2pt,topsep=2pt]\item $\vec v_r=(2,-1)$, $\vec v_s=(-4,2)$: paral·lels. $(1,2)\in s$? $\frac{1-3}{-4}=\frac12$, $\frac{2-1}{2}=\frac12$: sí, $\res{\text{coincidents}}$\item $\vec v_r=(1,-3)$, $\vec v_s=(2,-6)$: paral·lels. $(2,1)\in s$? $2=2\mu\Rightarrow\mu=1$, $1=4-6\mu\Rightarrow\mu=\frac12$: no, $\res{\text{paral·leles}}$\item $\vec v_r=(2,1)$, $\vec v_s=(1,1)$: secants. Substituïm $r$ en $s$: $-1+2\lambda-3-\lambda+1=0\Rightarrow\lambda=3$: $\res{(5,6)}$\item $\vec v_r=(-1,2)$, $\vec v_s=(1,-1)$: secants. $2(1+\lambda)+(1-\lambda)-4=0\Rightarrow\lambda=1$: $\res{(2,0)}$\end{enumerate}

\ex{2} $r:3x-3=2y+4\Rightarrow 3x-2y-7=0$; \ $s:y-1=-x+4\Rightarrow x+y-5=0$. \ $y=5-x\Rightarrow 3x-10+2x-7=0\Rightarrow x=\frac{17}{5}$, $y=\frac85$: $\res{\left(\tfrac{17}{5},\tfrac85\right)}$

\ex{2} \begin{enumerate}[label=\alph*),leftmargin=*,itemsep=2pt,topsep=2pt]\item $m_s=\frac{-6}{3}=-2$, $m_r=a$: $\res{a=-2}$ (i $r\neq s$: $(1,-2)\notin r$)\item $\frac36=\frac{b}{-4}\Rightarrow \res{b=-2}$. No poden ser coincidents: $\frac{-2}{1}\neq\frac12$.\end{enumerate}

\ex{3} \begin{enumerate}[label=\alph*),leftmargin=*,itemsep=2pt,topsep=2pt]\item $\frac a2=\frac{3}{-1}=\frac b4\Rightarrow \res{a=-6},\ \res{b=-12}$\item Punt de tall de les dues primeres: sumant, $3x-3=0$, $(1,2)$. Substituint en la tercera: $m+4-5=0\Rightarrow \res{m=1}$\end{enumerate}


% =====================================================
\seccio{8}{Rectes paral·leles i perpendiculars. Mediatriu}

\exf{1} $m=\frac12$: $y-4=\frac12(x-4)\Rightarrow \res{y=\tfrac12x+2}$. \ $r:y=3x-3$: $3x-3=\frac12x+2\Rightarrow x=2$, $y=3$: $\res{(2,3)}$

\ex{1} $x+3y+K=0$; \ $0+3+K=0\Rightarrow K=-3$: \ $\res{x+3y-3=0}$

\ex{1} $2x-3y+K=0$; \ $-2-6+K=0\Rightarrow K=8$: \ $\res{2x-3y+8=0}$

\ex{2} $s:3x+y+K=0$, $6-1+K=0$: $\res{3x+y-5=0}$. \ $t:x-3y+K=0$, $2+3+K=0$: $\res{x-3y-5=0}$. \ $y=2-3x\Rightarrow x-6+9x-5=0\Rightarrow x=\frac{11}{10}$, $y=-\frac{13}{10}$: $\res{\left(\tfrac{11}{10},-\tfrac{13}{10}\right)}$

\ex{1} $M(1,3)$, $\overrightarrow{AB}=(-6,6)\parallel(1,-1)$ (vector normal): $x-y+K=0$, $1-3+K=0$: $\res{x-y+2=0}$

\ex{1} \begin{tasks}(1)
  \task paral·lela: $\res{2x+5y-11=0}$; \ perpendicular: $\res{5x-2y+16=0}$
  \task paral·lela: $\res{x-4y+14=0}$; \ perpendicular: $\res{4x+y+5=0}$
  \task paral·lela: $\res{3x-y+9=0}$; \ perpendicular: $\res{x+3y-7=0}$
  \task paral·lela: $\res{3x+y+3=0}$; \ perpendicular: $\res{x-3y+11=0}$
\end{tasks}

\ex{1} \begin{enumerate}[label=\alph*),leftmargin=*,itemsep=2pt,topsep=2pt]\item $m=-3$: $y+4=-3(x-2)\Rightarrow \res{y=-3x+2}$\item $m'=-\frac14$: $y-1=-\frac14(x-4)\Rightarrow \res{y=-\tfrac14x+2}$\end{enumerate}

\ex{2} \begin{enumerate}[label=\alph*),leftmargin=*,itemsep=2pt,topsep=2pt]\item Mateix vector director $(2,-5)$: $\res{\left\{\begin{array}{l}x=1+2t\\ y=4-5t\end{array}\right.}$\item Director perpendicular a $(4,-1)$: $(1,4)$. \ $\res{\dfrac{x}{1}=\dfrac{y+3}{4}}$\end{enumerate}

\ex{2} \begin{enumerate}[label=\alph*),leftmargin=*,itemsep=2pt,topsep=2pt]\item $2x-1=5-x\Rightarrow x=2$, $y=3$; \ director $(3,-2)$: \ $\res{\dfrac{x-2}{3}=\dfrac{y-3}{-2}}$\item Sumant: $3x-6=0\Rightarrow(2,3)$; director perpendicular a $(2,-1)$: $(1,2)$: \ $\res{(x,y)=(2,3)+t(1,2)}$\end{enumerate}

\ex{2} \begin{enumerate}[label=\alph*),leftmargin=*,itemsep=2pt,topsep=2pt]\item $y+1=-\frac23(x-3)\Rightarrow \res{y=-\tfrac23x+1}$\item $m'=\frac32$: \ $\res{y-0=\tfrac32(x-0)}$, és a dir $y=\frac32x$\item Vector director $(3,-2)$: \ $\res{\dfrac{x}{3}=\dfrac{y-5}{-2}}$\item $y=-4$ és horitzontal: paral·lela $\res{y=2}$, perpendicular $\res{x=3}$\end{enumerate}

\ex{2} $M(1,-1)$, $\overrightarrow{AB}=(8,-4)\parallel(2,-1)$: $2x-y+K=0$, $2+1+K=0$: $\res{2x-y-3=0}$. \ $6-3-3=0$\checkmark; \ $d(P,A)=\sqrt{36+4}=\sqrt{40}=d(P,B)=\sqrt{4+36}$\checkmark

\ex{3} $\overrightarrow{AB}=(6,2)\parallel(3,1)$. Altura: $3x+y+K=0$, $9+7+K=0$: $\res{3x+y-16=0}$. \ Recta $AB$: $x-3y+2=0$. \ $x=3y-2\Rightarrow 10y-22=0$: $\res{H\left(\tfrac{23}{5},\tfrac{11}{5}\right)}$. \ $h=|\overrightarrow{HC}|=\sqrt{\tfrac{64}{25}+\tfrac{576}{25}}=\tfrac{8\sqrt{10}}{5}$, \ $|\overrightarrow{AB}|=2\sqrt{10}$, \ $S=\frac12\cdot2\sqrt{10}\cdot\frac{8\sqrt{10}}{5}=\res{16\ \text{u}^2}$

\ex{3} $\overrightarrow{AB}=(4,2)\parallel(2,1)$; perpendicular per $B$: $2x+y-8=0$; \ $x=1\Rightarrow y=6$: $\res{C(1,6)}$. \ $D=A+\overrightarrow{BC}=(-1,0)+(-2,4)=\res{D(-3,4)}$. \ $|\overrightarrow{AB}|=|\overrightarrow{BC}|=\sqrt{20}$: costats consecutius iguals i perpendiculars, és un quadrat (d'àrea 20 u²).


\end{document}
